% ============================================================
% pretty-charts 字体层：无衬线搭配（黑体 + Arial）
% 用法：\input{.../references/assets/tikz/fonts-sans.tex}
%
% 逐级回退：
%   Windows  → Arial + Microsoft YaHei
%   Linux    → TeX Gyre Heros + Noto Sans CJK SC（apt install fonts-noto-cjk）
%   兜底     → TeX Gyre Heros + FandolHei（TeX Live 自带）
%   再兜底   → 只发警告、不设字体，保证编译不因缺字体而失败
% 注意：Microsoft YaHei 是静态字体，可安全嵌入 PDF；Noto Sans SC 在 Windows 上多为
% 可变字体，xdvipdfmx 无法嵌入（见 spec/type.md §3），故这里不做该回退。
% ============================================================
\IfFontExistsTF{Arial}
  {\setmainfont{Arial}}                    % Windows / macOS
  {\IfFontExistsTF{TeX Gyre Heros}
     {\setmainfont{TeX Gyre Heros}}         % Linux 需 apt install fonts-texgyre
     {\PackageWarning{pretty-charts}{未找到西文无衬线字体（Arial / TeX Gyre %
       Heros），改用 XeLaTeX 默认西文字体；如需一致字形请安装 fonts-texgyre}}}
\IfFontExistsTF{Microsoft YaHei}
  {\setCJKmainfont{Microsoft YaHei}}       % 中文黑体（Windows，静态字体）
  {\IfFontExistsTF{Noto Sans CJK SC}
     {\setCJKmainfont{Noto Sans CJK SC}}    % Linux（fonts-noto-cjk）
     {\IfFontExistsTF{FandolHei}
        {\setCJKmainfont{FandolHei}}         % 兜底：TeX Live 自带
        {\PackageWarning{pretty-charts}{未找到中文黑体（已尝试 Microsoft YaHei / %
          Noto Sans CJK SC / FandolHei）；中文可能显示异常。请安装 fonts-noto-cjk}}}}
